% font: footnotesize
% TAB-07: generated by harness/analysis/paper_outputs.py; do not edit (every cell's source is in TAB-07.cells.json)
\begin{tabularx}{\linewidth}{>{\hsize=0.899\hsize\raggedright\hyphenpenalty=10000\exhyphenpenalty=10000\arraybackslash}X>{\hsize=1.784\hsize\raggedright\hyphenpenalty=10000\exhyphenpenalty=10000\arraybackslash}Xr>{\hsize=0.800\hsize\raggedleft\hyphenpenalty=10000\exhyphenpenalty=10000\arraybackslash}X>{\hsize=0.448\hsize\raggedleft\hyphenpenalty=10000\exhyphenpenalty=10000\arraybackslash}X>{\hsize=0.817\hsize\raggedright\hyphenpenalty=10000\exhyphenpenalty=10000\arraybackslash}X>{\hsize=1.383\hsize\raggedright\hyphenpenalty=10000\exhyphenpenalty=10000\arraybackslash}X>{\hsize=0.868\hsize\raggedright\hyphenpenalty=10000\exhyphenpenalty=10000\arraybackslash}Xl}
\toprule
Check & Quantity & Value & Interval & Interval level (\%) & Count & Rule & Stamped reading & Fired \\
\midrule
Fairness check & Opus: house-rule slips missed, of steps that should fire (counts) & 0 of 230 & --- & --- & --- & --- &  & --- \\
\addlinespace
Low-information flag & Haiku-medium single-sample accuracy at 6\ensuremath{\times}1 (pieces), mean over puzzles (\%) & 51.00 & --- & --- & --- & fires when that accuracy is above 60.0\% (as the analysis code applies it) &  & no \\
 & Median over puzzles of the same accuracy (\%) & 50.00 & --- & --- & --- & --- &  & --- \\
\addlinespace
Near-ceiling label & Opus-medium pooled accuracy, main run (\%) (label shown beside Test~1, Test~2) & 98.13 & \mbox{[95.63, 100.00]} & 95 & 80~puzzles & set on the length sweep at 6 steps: Opus-medium right on 20 of 20 chains and 20 of 20 pieces & near ceiling (Amendment~1) & yes \\
\addlinespace
Rewording check & Reworded prompt minus Haiku-medium on the same puzzles, both shapes (points) & +8.75 & \mbox{[0.00, 17.25]} & 95 & 20~puzzles in each shape & --- &  & --- \\
 & Reworded prompt minus Haiku-medium on the same puzzles, 1\ensuremath{\times}6 (points) & +11.50 & \mbox{[0.00, 23.50]} & 95 & 20~puzzles & --- &  & --- \\
 & Reworded prompt minus Haiku-medium on the same puzzles, 6\ensuremath{\times}1 (points) & +6.00 & \mbox{[\ensuremath{-}5.00, 17.00]} & 95 & 20~puzzles & --- &  & --- \\
\addlinespace
Drift check (top-up) & Top-up minus main run, Haiku-medium, both shapes (points) & +1.50 & \mbox{[\ensuremath{-}8.00, 10.00]} & 95 & 20~puzzles in each shape & --- &  & --- \\
 & Top-up minus main run, Haiku-medium, 1\ensuremath{\times}6 (points) & +6.00 & \mbox{[\ensuremath{-}3.00, 15.00]} & 95 & 20~puzzles & --- &  & --- \\
 & Top-up minus main run, Haiku-medium, 6\ensuremath{\times}1 (points) & \ensuremath{-}3.00 & \mbox{[\ensuremath{-}14.00, 8.00]} & 95 & 20~puzzles & --- &  & --- \\
 & Mean output tokens, relative change top-up against main (\%) & +0.29 & --- & --- & --- & --- &  & no \\
 & Accuracy condition & --- & --- & --- & --- & --- &  & no \\
 & Drift-check flag & --- & --- & --- & --- & fires if the pooled 95\% interval excludes 0 and \ensuremath{|}$d$\ensuremath{|}~\ensuremath{>}~10 points, or mean output tokens move by more than 25\%; a flag, not a gate & not fired & no \\
\bottomrule
\end{tabularx}
% note: Checks and flags, not tests. Each Quantity cell names its unit: \% for shares and accuracies, points for differences, ratios and counts as named. Dashes mark cells with no stamped value for that column.
% note: Terms. A slip is a step whose written balances differ from that step's rule applied to the model's own previous balances. T steps are steps of the targeted rule kinds: kinds added, or made more frequent, where Haiku slipped while the puzzles were tuned; G steps are steps of the other kinds. Same-slip pairwise agreement (Flag~R): among pairs of Haiku-medium samples that slip at the same step of the same puzzle (on a chain only its first slip counts, on pieces every slipped piece), the share that wrote the same wrong balances there. Slip-rate ratio: a model's slip rate on T steps divided by its slip rate on G steps. House variant: the house rule is a puzzle-wide rule that should act on some steps; the ratio's numerator is the share of those steps (the exposed steps) where the model missed it (the exposed slips). Targeting index: Haiku-medium's ratio divided by the comparison's ratio, here Opus-medium and Sonnet-medium pooled; near one means the T kinds are no harder for Haiku in particular. Zero-slip corrected: the model had no slips in one group, so half a slip was added to both groups' counts before dividing, which pulls the ratio towards one.
% note: Rule column: the rule text is stamped for the flag of the drift check between runs and for the near-ceiling label's source (given in plain words), and the format flag's rule is stamped and given under the format panels. For Flag~R, the targeting indices behind Flag~K and the low-information flag, the result file stamps no threshold: the Rule cell gives each threshold as the Study~1 analysis code that produced the result file applies it (the fairness check's code and the analysis driver, at the code commit ⟦TAB-23⟧ lists for the Study~1 analysis). The thresholds printed in those Rule cells are the ones the pre-registered plan (Flag~R and Flag~K) and Study~1's reporting definitions (the low-information flag) state. The stamped reading column repeats the files' own reason strings, in plain words where they hold code names.
% note: Assessability: the fewest slips and steps an index needs to be assessable are as the analysis code sets them. For the house variant it asks the comparison (Opus and Sonnet pooled) for at least 3 exposed slips, the number in its reading; the pre-registered plan states no such number, and applies its general rule to the house index, which asks the comparison for at least 5 slips in each group, with its house-rule slips as one group. With 0 exposed slips, the house variant is not assessable under either number.
% note: The rewording check's lower interval ends at zero are floating-point zeros in the file (exact zeros up to rounding error), printed as zero without a sign.
% note: The rewording check (a robustness check planned in the pre-registered design, outside the confirmatory tests; no test): its wording, puzzles and comparison were fixed after Lock~3, before any result was opened. It was checked on Haiku-medium single samples only.
% note: Near ceiling (Amendment~1): a label that the pre-registration's Amendment~1 turned on for Tests~1 and 2 because their comparator, Opus-medium, was right on 40 of 40 answers in the length sweep at 6 steps, a pooled whole-answer accuracy of 100.00\%, at or above the pre-set 95\%; it marks a comparator at or near a perfect score. With it, the error-rate ratio (the challenger's errors over Opus-medium's) is reported with its interval next to the no-worse-than reading (in ⟦TAB-02⟧); where Opus-medium made no errors the ratio is undefined, and its exact 95\% upper error rate, in \%, is given instead.
% note: The near-ceiling label was set from the length sweep (Amendment~1) and was not re-decided afterwards; the main run's Opus-medium pooled score (the near-ceiling row) is shown for information only.
% note: The fairness check's first-slip variant is a separate stamped sensitivity and is not shown here.
